\section{Seventeen integers: compressing congruence witnesses}
\label{sec:integer17}

The following constructions use unrestricted integer coordinates. They give
the variable--degree pairs $(23,10)$, $(17,88)$, and $(16,184)$,
with a $(16,192)$ sum-of-squares control.
The first two improve the incoming $(24,10)$ and $(19,88)$ points
\cite{budget,joint}; the last trades degree for another coordinate.
Their literal circuits use respectively $131$, $114$, and $135$ binary
operations; the degree-192 control uses $134$. Scalar arithmetic and
squaring are counted. These are different
tradeoffs from the much higher-degree ten-integer reductions. In particular,
the $114$-operation circuit is larger than the incoming $109$-operation
nineteen-integer circuit. No global minimum or historical priority is claimed.

All four use the same six fixed integers
\[
 \beta=2^{1064921}-1,\qquad K_0,M_C,M_F,\qquad
 A_c=2^{1113},\quad B_c=2^{666104}.
\]
The first four are the entire source-34
\texttt{mixed-parity-unrounded-1064921} tuple; $M_F$ already includes its
shift by $\beta$. The last two are source71's determinant scales
\cite{source71}. The actual six integers, compressed minimal little-endian
bytes, finite native recipe, and complete raw hashes are supplied in the
integer component's \texttt{coefficients/} directory. The raw encodings
total $615\,883$ bytes. These fixed leaves are not additional unknowns.
They satisfy
\[
 \beta\ge17,\quad\beta\equiv1\pmod3,\quad
 0<M_C<\beta<M_F<2\beta,\quad K_0\ge0,\quad A_c,B_c\ge2.
\]

\subsection{Two inherited guards}

Put $\operatorname{Pos}(p_1,p_2,p_3)=1+p_1^2+p_2^2+p_3(p_3+1)$.
It takes exactly the positive integer values, by the three-square theorem
applied to $4n-3$. The nearest-square joint-positivity lemma used in the
incoming constructions is
\begin{equation}\label{eq:integer-pair}
 x,y>0\quad\Longleftrightarrow\quad
 \exists t,p_1,p_2,p_3\in\Z:\quad
 xy-(t^2-x-y+1)^2=\operatorname{Pos}(p_1,p_2,p_3).
\end{equation}
Indeed a positive left side forces $xy>0$ and excludes both inputs being
negative. Conversely, for $x,y>0$, choose the nearer neighboring square
to $x+y-1$. Its squared error is less than $x+y-1\le xy$, leaving a
positive integer represented by $\operatorname{Pos}$.

For $\alpha=q-F-C$ define
\[
 d=(A_cF)(B_cC)-(q-A_cF-B_cC)^2.
\]
The inequality $d>0$ forces $F,C,\alpha,q$ to have one common nonzero
sign. To see this, put $U=A_cF$, $B=B_cC$. Then $U,B$ have a common
sign and $|q-U-B|<\sqrt{UB}\le(|U|+|B|)/2$. Hence $q$ has their sign
and $|q|>(|U|+|B|)/2\ge|F|+|C|$, which also fixes the sign of
$\alpha$. Source71's separate canonical-witness estimates prove
completeness for this determinant; that calibrated existence restriction
is inherited, not re-proved by finite sign tests.

\subsection{The complete polynomial}

The seventeen independent integer coordinates are
\[
 (J,C,z,f,h,r,v,w,\tau,\nu,\gamma,T,c,t,p_1,p_2,p_3).
\]
The following are polynomial abbreviations only:
\begin{align*}
 J_0&=3J+1,&q&=\beta J_0+1,&W&=4(3w+1),\\
 X&=qW,&Y&=(2v+1)q^3,&E&=XY,\\
 a&=(X+1)Y,&H&=4a+3,&\Delta&=a^2+H,\\
 F&=(K_0+W)C-(q-1)z,&\alpha&=q-F-C.&&
\end{align*}
Continue with
\begin{align*}
 R&=(q-F)(q^3-q)-C(q^2-1)+(M_C+qM_F)J_0,\\
 k&=R+1+hE,\qquad \eta=c-Yk,\qquad \zeta=k-\eta,\\
 I&=(2r+1)c^2,\qquad V=c(Tf-1)-Rf^2,\\
 m&=XY^2,\qquad B_0=mk,\qquad s_0=X+\gamma H,\qquad Z=\eta\zeta.
\end{align*}
Define five residuals:
\begin{align}
 A_1&=\tau^2+(2\tau-k)B_0-1,\nonumber\\
 A_2&=2acs_0+s_0^2-Hc^2-1,\nonumber\\
 A_3&=I^2-\Delta(f^2-1),\label{eq:integer17-rows}\\
 A_4&=(V+\nu)^2-I^2\nu(2V+\nu)-1,\nonumber\\
 A_5&=dZ-(t^2-d-Z+1)^2-\operatorname{Pos}(p_1,p_2,p_3).\nonumber
\end{align}
The ordinary integer polynomial is
\begin{equation}\label{eq:integer17}
 \boxed{P_{17}=A_1^2+A_2^2+A_3^2+A_4^2+A_5^2.}
\end{equation}

The centered roots $\tau_{\rm old}=\tau+B_0$ and
$y_{\rm old}=\nu+V$ give the identities
\[
 A_1=\tau_{\rm old}^2-m(m+1)k^2-1,\qquad
 A_4=I^2(V^2-y_{\rm old}^2)+y_{\rm old}^2-1.
\]
Also $D=ac+s_0$ gives $A_2=D^2-\Delta c^2-1$.

\subsection{One quotient replaces two}

\begin{lemma}[Integer quotient compression]\label{lem:integer-quotient}
Suppose $c\ne0$ and $c\mid f^2-1$. The simultaneous integer identities
\[
 V=of-c=jc-R
\]
are equivalent, after projection of $j,o$, to
$V=c(Tf-1)-Rf^2$ for one integer $T$.
\end{lemma}
\begin{proof}
From the old identities, $c\mid of+R$. The identity
\[
 f(of+R)-o(f^2-1)=o+Rf
\]
shows that $T=(o+Rf)/c$ is integral and has the desired value of $V$.
Conversely put $b=(f^2-1)/c$ and define
\[
 o=cT-Rf,\qquad j=Tf-1-Rb.
\]
These are integers and restore both identities. The witness-restoration
division introduces neither division nor a $b$ coordinate into $P_{17}$.
\end{proof}

The order in which $c\mid f^2-1$ is established matters. Invoking an
auxiliary-index theorem that already used the deleted congruences would
be circular. Here the needed bootstrap is source70's signed orientation
argument \emph{before} its auxiliary-congruence step
\cite{source70}. We spell out its dependency boundary.

At a $P_{17}$ zero, $A_5=0$ gives $d>0$ and $\eta\zeta>0$.
Write $Q=|q|$, $U=|X|$, $Z_0=|Y|$, $K=|k|$, and $M=|E|$.
The determinant, authentic mask bounds, and affine residue restrictions give
\[
 Q\ge16,\quad Q(Q^2-1)<R<Q^4-2,\quad
 U\ge4Q,\quad Z_0\ge Q^3,\quad M>3R+1.
\]
The positive product gives $K\ge2$, $0<|\eta|<K$, and
$K(Z_0-1)<|c|<K(Z_0+1)$. Let $B=|a+2|$ and $P=|2XY^2+1|$.
They satisfy $P>B\ge2$ and $B>R$. Classifying the first and main
Pell norms gives
\[
 K=2\psi_P(n_0),\qquad |c|=\psi_B(p),\qquad p>n_0.
\]
The recurrence modulo $M$ and $k=R+1+hE$ give
$\pm2n_0\equiv R+1\pmod M$. The positive branch gives
$n_0\ge(R+1)/2$, while the negative branch gives $n_0>R$.
Thus $p>n_0\ge24$; Pell growth yields
\[
 |c|>B(B^2-1)^2,\qquad |c|>2p,\qquad |c|>2R.
\]
The inherited ordinary strong-rank lemma can now be applied to $A_3$.
It gives $|f|=\chi_B(l)$ with $p\mid l$, without using either missing
congruence. Therefore $\psi_B(p)\mid\psi_B(l)$ and
\[
 c\mid (B^2-1)\psi_B(l)^2=f^2-1.
\]
The precise retained source proof is
\texttt{upstream/ORIENTATION\_PROOF.md}, Sections 3--4.2; its
auxiliary-congruence argument starts in Section 4.3.

Lemma~\ref{lem:integer-quotient} restores $j,o$. All other deleted
coordinates are monic polynomial definitions, the two root shifts are
integral, and the old positive guard triples can be restored by
$\operatorname{Pos}$. This gives a zero of the incoming calibrated
37-integer affine quartic \cite{round11}. Conversely, any zero of that
source supplies the same divisibility, hence $T$; its transport row gives
$F=(K_0+W)C-(q-1)z$ after the invertible change
$\alpha=q-F-C$. Forget the monic coordinates, perform the root shifts,
and use \eqref{eq:integer-pair}. Thus $P_{17}$ has an integer zero
exactly when that source does. This is a new existence equivalence on
its full integer solution set, with no additional canonical restriction.

The exact verifier restores all 22 arithmetic rows of the literal parent:
four match the new arithmetic residuals, one is $\pm R(cb-f^2+1)$,
and all others vanish identically. It thereby exposes the sole use of the
proved divisibility rather than silently treating $b$ as an unknown.

\section{Sixteen integers: merging an ordinary and a rational Pell root}
\label{sec:integer16}

\begin{lemma}[Separated rational norm]\label{lem:integer-rational-norm}
Let $A,Q,S,V\in\Z$, with $A>0$, $Q>1$, and $S^2>A$.
Put $\delta=Q-1$ and $B_{\rm num}=S^2(QV^2-1)$. Then
\[
 \exists n\in\Z:\quad
 [\delta(n^2-A)-B_{\rm num}]^2-4\delta A B_{\rm num}=0
\]
if and only if there are integers $u,y$ with
$u^2=A$ and $QV^2-(Q-1)y^2=1$.
\end{lemma}
\begin{proof}
If $V=0$, the displayed expression is a square plus the strictly positive
integer $4\delta AS^2$. Thus any zero has $|V|\ge1$. The rational number
$B=B_{\rm num}/\delta$ then satisfies $B\ge S^2>A$.
Dividing by $\delta^2$ gives $(n^2-A-B)^2=4AB$.
Here $n\ne0$, since $n=0$ would give $A=B$. The identity
\[
 (n^2+A-B)^2-4n^2A=(n^2-A-B)^2-4AB
\]
shows that $u=(n^2+A-B)/(2n)$ is a rational square root of the integer
$A$, hence an integer. The rational $s=n-u$ satisfies $s^2=B$, so
$y=s/S$ is rational and obeys $QV^2-(Q-1)y^2=1$.

Source69's adjacent-coefficient Pell descent makes this $y$ integral.
For completeness, if $y=a/d$ is reduced with $d>0$, the norm gives
$d^2\mid Q-1$. Make $V,y$ positive by changing signs. Unless $V=1$,
the integral coordinate
\[
 V'=(2Q-1)V-2(Q-1)y
\]
and rational $y'=(2Q-1)y-2QV$ preserve the norm. The reduced denominator
of $y'$ remains $d$, since $2Q-1\equiv1\pmod d$.
The inequalities $y>V$ and $QV>(Q-1)y$ imply $-V<V'<V$; the norm
excludes $V'=0$. Strong induction on $|V|$ reaches $V=1$, where
$y=1$ and $d=1$. This is the inherited descent proof, applied to the
new separated norm. Conversely, take $n=u+Sy$ and factor the norm.
\end{proof}

Delete $\tau,\nu$ from the seventeen coordinates and add one unrestricted
integer $n$. Keep every other abbreviation above and put
\begin{align*}
 A_0&=1+m(m+1)k^2,\qquad S=k(2m+1),\qquad\delta=I^2-1,\\
 L&=\delta\bigl[n^2+2n(B_0+SV)+2B_0SV-mk^2-1\bigr]
       +S^2(1-V^2),\\
 N&=L^2-4\delta A_0S^2(I^2V^2-1).
\end{align*}
The new equation is the four-square polynomial
\begin{equation}\label{eq:integer16}
 \boxed{P_{16}=A_2^2+A_3^2+A_5^2+N^2.}
\end{equation}
It is exactly the norm of Lemma~\ref{lem:integer-rational-norm} with
$Q=I^2$ and $n_{\rm total}=n+B_0+SV$, because
\[
 \delta(n_{\rm total}^2-A_0)-S^2(I^2V^2-1)=L.
\]
This centering cancels the two leading $V^2$ terms and lowers the degree.

\paragraph{The separation does not assume either missing norm.}
From $A_5=0$ alone, $d>0$ and $\eta\zeta>0$. The determinant and
fixed coefficient bounds imply $|q|\ge16$; the definitions give
$|Y|\ge|q|^3$ and $|X|\ge4|q|$. As $\eta,\zeta$ have the same
sign, $|k|\ge2$ and
\[
 |c|>|k|(|Y|-1)>1.
\]
Hence $I^2>1$, while $|m|>1$ gives $A_0>0$. Finally the exact identity
\begin{equation}\label{eq:integer-separation}
 S^2-4A_0=k^2-4\ge0
\end{equation}
proves $S^2>A_0$. All hypotheses of the rational-norm lemma are therefore
established before either deleted Pell root is extracted.

At a zero of $P_{16}$, the lemma supplies integral $\tau_{\rm old}$
and $y_{\rm old}$; put $\tau=\tau_{\rm old}-B_0$ and
$\nu=y_{\rm old}-V$. This restores a zero of $P_{17}$. Conversely a
$P_{17}$ zero gives a $P_{16}$ zero by $n=\tau+S\nu$.
This proof covers arbitrary signed coordinates, and introduces no new
divisibility restriction on canonical witnesses.

\subsection{Exact degree certificates}

For $P_{17}$ the degrees of $X,Y,E,a,R,k$ are $2,4,6,6,5,7$;
$\eta,\zeta$ have degree 11 and $d$ has degree 4.
The five residual bounds are $24,14,14,14,44$.
Only the joint guard reaches degree 44. Its top part is
$-\eta_{\rm top}^4$, with
\[
 \eta_{\rm top}=-2^4 3^8\beta^7hJ^7wv^2.
\]
Consequently
\[
 [h^8J^{56}w^8v^{16}]P_{17}=2^{32}3^{64}\beta^{56}>0,
\]
proving exact degree 88 on the actual fixed coefficient slice.

For $P_{16}$, $\deg S=17$ and $\deg V=7$.
The top part of $L$ is $-S_{\rm top}^2V_{\rm top}^2$, of degree 48;
the second term of $N$ has degree at most 94. Thus $N$ has degree 96.
Specifically,
\begin{align*}
 S_{\rm top}&=2^8 3^{13}\beta^{11}hJ^{11}w^2v^3,\\
 V_{\rm top}&=f^2(324\beta^3J^3wC-81\beta^4J^4z),
\end{align*}
and
\[
 [h^8J^{120}w^{16}v^{24}z^8f^{16}]P_{16}
   =2^{64}3^{136}\beta^{120}>0.
\]
The other squared rows have degrees at most $28,28,88$, proving exact
degree 192. The direct uncentered weighted norm would have upper degree
216. Exact homogeneous-component arithmetic, rather than a naive gate
bound, verifies both coefficients in the supplied auditor.

\subsection{A degree-184 finalizer}

Let $G=A_2^2+A_3^2+A_5^2$. Instead of $G+N^2$, define
\begin{equation}\label{eq:integer16unit}
 \boxed{P_{16}^{\rm unit}=G+N+GN=(1+G)(1+N)-1.}
\end{equation}
For every integer assignment $G\ge0$. If the new polynomial is zero,
the positive integer $1+G$ multiplies the integer $1+N$ to give 1;
therefore $G=N=0$. Conversely these equalities make both polynomials
zero. This exact zero-set equivalence needs no sign or congruence
assumption about $N$.

The degree is $88+96=184$, certified by
\[
 [h^{12}J^{116}w^{16}v^{28}z^4f^8]P_{16}^{\rm unit}
   =2^{64}3^{132}\beta^{116}>0.
\]
The expanded form $G+N+GN$ uses $135$ binary operations, one more than
the degree-192 control. The archive retains both explicit polynomials.

\section{A lower-degree companion in twenty-three integers}
\label{sec:integer23}

This companion combines the incoming residue orientation with the paired
Pell-root translations \cite{round11,budget}. Its independent coordinates are
\[
 \begin{gathered}
 J,F,\alpha,z,f,h,r,v,w,\tau,\eta,\zeta,\nu,\gamma,o,j,\ E,a,I,\ t,p_1,p_2,p_3.
 \end{gathered}
\]
Here $F,\alpha,E,a,I$ are coordinates, not their expressions in the
seventeen-variable equation. Define
\begin{align*}
 J_0&=3J+1,&q&=\beta J_0+1,&W&=4(3w+1),\\
 X&=qW,&Y&=a-E,&k&=\eta+\zeta,\\
 c&=Yk+\eta,&V&=fo-c,&C&=q-F-\alpha,\\
 H&=4a+3,&\Delta&=a^2+H,&B_0&=E(c-\eta),\\
 s_0&=X+\gamma H.&&&&
\end{align*}
Use
\[
 R=((q-1)(q-F)+\alpha)(q^2-1)+(M_C+qM_F)J_0
\]
and the earlier definitions of $d$, $Z=\eta\zeta$, and $A_5$.
The eleven residuals, whose squared sum defines $P_{23}$, are
\begin{align*}
 &a-E-(2v+1)q^3, && E-X(a-E),\\
 &\tau^2+(2\tau-k)B_0-1,&&2acs_0+s_0^2-Hc^2-1,\\
 &k-1-hE-R,&&(K_0+W)C-(q-1)z-F,\\
 &V+k-1-cj-hE,&&I-(2r+1)c^2,\\
 &I^2-\Delta(f^2-1),&&(V+\nu)^2-I^2\nu(2V+\nu)-1,\\
 &A_5.&&
\end{align*}
The count is $37-6+4-12=23$: replace the source affine quartic's six
private guard coordinates by four, translate both roots, and delete twelve
monic coordinates. Their exact restorations are
\begin{align*}
 N&=q^3,&X&=qW,&L&=q^2-1,&C_2&=c^2,\\
 F_2&=f^2,&B_{\rm aux}&=I^2,&g&=\nu(2V+\nu),&\Delta&=a^2+4a+3,\\
 D&=ac+s_0,&P&=k+B_0,&c&=(a-E)k+\eta,&V&=fo-c,
\end{align*}
with $\tau_{\rm old}=\tau+B_0$, $y_{\rm old}=V+\nu$.
Every arithmetic map is integral and polynomial. The auditor checks all
22 source arithmetic rows: twelve vanish identically and ten match the
displayed arithmetic rows once each, up to sign. It also authenticates
both old guard arguments.

The residual degrees are at most
$4,3,4,5,4,2,3,5,4,5,4$. Only the $I$ definition contains $r$, so
\[
 [r^2a^4\eta^4]P_{23}=4.
\]
Thus the degree is exactly ten, independently of the fixed ports.

\begin{theorem}[Inherited calibration, new integer reductions]
For each $P\in\{P_{23},P_{17},P_{16},P_{16}^{\rm unit}\}$, composition of the new
reductions above with the named source native/compiler, Pell, and
syntactic-refutation calibration theorems proves
\[
 \mathsf{ZFC}\ \vdash\
  \bigl(\neg\exists\mathbf x\in\Z^{\dim P}:P(\mathbf x)=0\bigr)
  \ \longleftrightarrow\ \ConZFC.
\]
\end{theorem}

The theorem asserts the equivalence, not either side. The full native and
logical calibration is inherited; the finite algebra/data checks are not an
end-to-end Lean/Rocq or encoded internal-ZFC certificate. All four literal
packets, inlined expressions, actual coefficients, exact degree receipts,
source restorations, and shared DAGs are included. The verifier uses only
Python's standard library and works with assertions disabled. This component's
new steps have independent mathematical review within the AI collaboration;
that is not external human peer review.
